\begin{afterword}
\summaryhead

The post argues that the foolishness or wickedness of a claim's believers is little evidence against the claim. Flying saucer cults would arise whether or not aliens had visited, so their existence says little either way. As Robert Pirsig put it, the world's greatest fool saying the Sun is shining does not make it dark. Someone wrong 99.99\% of the time could be reversed into a guide that is right 99.99\% of the time, but to be that wrong they would need to know the truth.

The same holds for evil people. The ``horns effect'' leads people to treat everything Stalin said as false, yet Stalin also believed that 2 + 2 = 4. The post then lists corollaries: argue against an idea's strongest advocates; lunatics who hold an idea are no evidence against it; ad hominem is not valid; an argument against genocide should not depend on Hitler; willingness to believe follows willingness to affiliate; and reversing a failed approach, in computers or in artificial intelligence, is unlikely to find the one that works.

\responsehead

The post's central claim is correct and well argued. That fools or villains believe something tells you little about whether it is true, because they are not reliably wrong. The flying saucer case applies the definition of evidence properly and names the one assumption that would change the answer. The 99.99\% example shows why reliable error would take knowledge. The point is old, and the post's own source says so: Pirsig offers the line about the fool and the sun to illustrate ``an old rule of logic.'' The first corollary, argue against the strongest advocates, goes back at least to Mill's \textsc{On Liberty}.

The step from stupidity to evil is asserted, not argued. ``How much less should human evil anticorrelate with truth?'' has no reason behind its ``how much less.'' By the post's own account, being reliably wrong takes someone who knows the truth, and a deliberate liar does. The Stalin example supports a narrower claim, that an evil person is not wrong about everything. The narrower claim is all the post needs.

The corollary ``Ad hominem argument is not valid'' also needs qualifying. It is true in the logician's sense: attacking the speaker does not show the claim false. As a rule about what to believe, it goes further than the post's principle. The principle says that a fool's word is weak evidence, not evidence of the opposite, and that leaves a speaker's reliability relevant whenever a claim depends on the speaker's word. ``Argument Screens Off Authority,'' posted two days later and placed near this post in the book, allows for this.

In short: a correct and old point about evidence, well illustrated, extended to evil people without an argument, and followed by a rule on ad hominem that holds in the logician's sense but not as a rule about whose word to trust.
\end{afterword}
